% !TEX program = xelatex
% !TEX encoding = UTF-8
%% ============================================================
%%  Analiță și Algebră pentru Olimpiadă — VOLUM COMPLET
%%  Teorie + Probleme de Analiză + Probleme de Algebră
%%  Compilare: xelatex (x3) pentru cuprins complet
%% ============================================================
\documentclass[11pt,a4paper]{report}

%% ---- encoding (compatibil pdflatex si xelatex) ----
\usepackage{iftex}
\ifXeTeX
  \usepackage{fontspec}
  \setmainfont{Latin Modern Roman}
\else
  \usepackage[utf8]{inputenc}
  \usepackage[T1]{fontenc}
\fi

%% ---- limbă ----
\usepackage[romanian]{babel}

%% ---- matematică ----
\usepackage{amsmath,amssymb,amsthm}
\usepackage{mathtools}
\usepackage{mathrsfs}

%% ---- layout ----
\usepackage{multicol}
\usepackage{graphicx}
\usepackage{caption}
%% KDP: lets TeX loosen a line instead of pushing inline math into the margin
\setlength{\emergencystretch}{3em}
%% KDP: a heading never stays alone at the bottom of a page. Otherwise a breakable
%% tcolorbox that cannot start under it is pushed to the next page and drawn over
%% the running head (tcolorbox: "Using nobreak failed"). Same trick as needspace.sty.
\makeatletter
\newcommand{\needlines}[1]{\par\begingroup\dimen@=#1\baselineskip
  \vskip\z@\@plus\dimen@ \penalty-100 \vskip\z@\@plus-\dimen@
  \vskip\dimen@ \penalty9999 \vskip-\dimen@ \vskip\z@skip\endgroup}
\makeatother
\AtBeginDocument{%
  \pretocmd{\section}{\needlines{9}}{}{\typeout{needlines: section not patched}}%
  \pretocmd{\subsection}{\needlines{8}}{}{\typeout{needlines: subsection not patched}}%
  \pretocmd{\subsubsection}{\needlines{6}}{}{\typeout{needlines: subsubsection not patched}}}
\usepackage{geometry}
%% Geometry identical to the printed books (VRAI LIVRES / interieurs_corriges).
%% Values measured from the shipped PDFs:
%%   - text block: 356.96pt = 4.9578in wide  => left = right = 2.2cm
%%   - header: top margin 34.64pt + headheight 22pt => head rule at 56.64pt
%%   - footer: page-number baseline 32.77pt above the bottom edge
%% Do NOT change these without regenerating the covers: the spine width depends
%% on the page count, and the page count depends on the text block.
\geometry{paperwidth=6.69in,paperheight=9.45in,
          left=2.2cm,right=2.2cm,top=34.64pt,bottom=32.77pt,
          includehead,includefoot,headsep=8pt,footskip=16pt}
\usepackage{enumitem}
\usepackage{xcolor}
\usepackage{tikz}
\usepackage{tcolorbox}
\tcbuselibrary{skins,breakable}
\usepackage{titlesec}
\usepackage{fancyhdr}
\usepackage[colorlinks=true,linkcolor=blue!60!black,urlcolor=blue!60!black]{hyperref}

%% ---- paleta cromatică unificată ----
\definecolor{ink}{HTML}{1C2230}
\definecolor{grey}{HTML}{6B7280}
\definecolor{olm}{HTML}{2A9D8F}
\definecolor{ojm}{HTML}{E76F51}
\definecolor{onm}{HTML}{6A4C93}
\definecolor{solcol}{HTML}{0B7A52}
\definecolor{teoblue}{RGB}{25,55,105}
\definecolor{obsgray}{RGB}{90,90,90}
\definecolor{acc}{RGB}{80,40,10}

%% ---- box-uri tehnice ----
\newtcolorbox{teobox}[1]{%
  enhanced,breakable,sharp corners,
  colback=onm!6,colframe=onm,
  colbacktitle=onm,coltitle=white,
  borderline west={4pt}{0pt}{onm},
  boxrule=0.4pt,left=10pt,top=4pt,bottom=4pt,
  title={#1},fonttitle=\bfseries}
\newtcolorbox{tehbox}[1]{%
  enhanced,breakable,sharp corners,
  colback=ojm!6,colframe=ojm,
  colbacktitle=ojm,coltitle=white,
  borderline west={4pt}{0pt}{ojm},
  boxrule=0.4pt,left=10pt,top=4pt,bottom=4pt,
  title={#1},fonttitle=\bfseries}

%% ---- teoreme ----
\theoremstyle{plain}
\newtheorem{teorema}{Teoremă}[chapter]
\newtheorem{propozitie}[teorema]{Propoziție}
\newtheorem{lema}[teorema]{Lemă}
\newtheorem{corolar}[teorema]{Corolar}
\theoremstyle{definition}
\newtheorem{definitie}[teorema]{Definiție}
\newtheorem{exemplu}[teorema]{Exemplu}
\newtheorem{aplicatie}[teorema]{Aplicație}
\theoremstyle{remark}
\newtheorem{observatie}[teorema]{Observație}
% obs redefinit mai jos ca newenvironment

%% ---- bare colorate pe teoreme ----
\tcolorboxenvironment{teorema}{%
  enhanced,breakable,sharp corners,boxrule=0pt,frame hidden,
  borderline west={4pt}{0pt}{onm},
  colback=onm!5,left=10pt,right=4pt,top=3pt,bottom=3pt}
\tcolorboxenvironment{propozitie}{%
  enhanced,breakable,sharp corners,boxrule=0pt,frame hidden,
  borderline west={3.5pt}{0pt}{onm!75},
  colback=onm!3,left=10pt,right=4pt,top=3pt,bottom=3pt}
\tcolorboxenvironment{lema}{%
  enhanced,breakable,sharp corners,boxrule=0pt,frame hidden,
  borderline west={3pt}{0pt}{onm!55},
  colback=onm!2,left=10pt,right=4pt,top=3pt,bottom=3pt}
\tcolorboxenvironment{corolar}{%
  enhanced,breakable,sharp corners,boxrule=0pt,frame hidden,
  borderline west={2.5pt}{0pt}{onm!40},
  colback=white,left=10pt,right=4pt,top=2pt,bottom=2pt}
\tcolorboxenvironment{definitie}{%
  enhanced,breakable,sharp corners,boxrule=0pt,frame hidden,
  borderline west={4pt}{0pt}{olm},
  colback=olm!5,left=10pt,right=4pt,top=3pt,bottom=3pt}
\tcolorboxenvironment{exemplu}{%
  enhanced,breakable,sharp corners,boxrule=0pt,frame hidden,
  borderline west={3.5pt}{0pt}{ojm},
  colback=ojm!4,left=10pt,right=4pt,top=3pt,bottom=3pt}
\tcolorboxenvironment{aplicatie}{%
  enhanced,breakable,sharp corners,boxrule=0pt,frame hidden,
  borderline west={3pt}{0pt}{ojm!80},
  colback=ojm!3,left=10pt,right=4pt,top=3pt,bottom=3pt}
\tcolorboxenvironment{observatie}{%
  enhanced,breakable,sharp corners,boxrule=0pt,frame hidden,
  borderline west={2.5pt}{0pt}{grey},
  colback=white,left=10pt,right=4pt,top=2pt,bottom=2pt}

%% ---- macro-uri algebră ----
\newcommand{\NN}{\mathbb{N}}
\newcommand{\ZZ}{\mathbb{Z}}
\newcommand{\QQ}{\mathbb{Q}}
\newcommand{\RR}{\mathbb{R}}
\newcommand{\CC}{\mathbb{C}}
\DeclareMathOperator{\ord}{ord}
\DeclareMathOperator{\Ker}{Ker}
\DeclareMathOperator{\Ima}{Im}
\DeclareMathOperator{\car}{car}
\DeclareMathOperator{\Hom}{Hom}
\DeclareMathOperator{\End}{End}
\DeclareMathOperator{\Aut}{Aut}
\DeclareMathOperator{\Inn}{Inn}
\DeclareMathOperator{\Orb}{Orb}
\DeclareMathOperator{\Stab}{Stab}
\newcommand{\cl}[1]{\widehat{#1}}
\DeclareMathOperator{\Tr}{Tr}
\DeclareMathOperator{\rang}{rang}
\DeclareMathOperator{\Span}{Span}
\DeclareMathOperator{\adj}{adj}
\DeclareMathOperator{\sgn}{sgn}
\DeclareMathOperator{\diag}{diag}
\DeclareMathOperator{\Spec}{Spec}
\newcommand{\GL}{\mathrm{GL}}

%% ---- macro-uri culegeri (titluri colorate) ----
\newcommand{\capitolcurent}{Introducere}
\newcommand{\sectiunealg}[2]{%
  \clearpage\vspace*{6mm}%
  \addcontentsline{toc}{section}{\textcolor{blue!60!black}{#1}}%
  {\noindent\color{ink}\rule{\linewidth}{3pt}}\par\vspace{2mm}%
  {\noindent\fontsize{34}{38}\selectfont\bfseries\color{ink}#1}\par\vspace{2mm}%
  {\noindent\large\itshape\textcolor{grey}{#2}}\par\vspace{2mm}%
  {\noindent\color{ink}\rule{\linewidth}{3pt}}\par\vspace{7mm}}
\newcommand{\parteamare}[2]{%
  \clearpage\vspace*{4mm}%
  \addcontentsline{toc}{subsection}{#1}%
  {\noindent\color{ink}\rule{\linewidth}{2.4pt}}\par\vspace{2mm}%
  {\noindent\fontsize{30}{34}\selectfont\bfseries\color{ink}#1}\par\vspace{1.5mm}%
  {\noindent\large\itshape\textcolor{grey}{#2}}\par\vspace{2mm}%
  {\noindent\color{ink}\rule{\linewidth}{2.4pt}}\par\vspace{6mm}}
\newcommand{\parteamareenunt}[1]{%
  \vspace*{4mm}%
  \addcontentsline{toc}{subsection}{#1}%
  {\noindent\color{ink}\rule{\linewidth}{2.4pt}}\par\vspace{2mm}%
  {\noindent\fontsize{30}{34}\selectfont\bfseries\color{ink}#1}\par\vspace{1.5mm}%
  {\noindent\color{ink}\rule{\linewidth}{2.4pt}}\par\vspace{6mm}}
\newcommand{\subsectiuneclasa}[3]{%
  \clearpage\renewcommand{\capitolcurent}{#3}\vspace*{3mm}%
  \addcontentsline{toc}{subsection}{\quad #1}%
  {\noindent\color{ink}\rule{\linewidth}{1.6pt}}\par\vspace{2mm}%
  {\noindent\Huge\bfseries\color{ink}#1}\par\vspace{1.5mm}%
  {\noindent\large\itshape\textcolor{grey}{#2}}\par\vspace{2mm}%
  {\noindent\color{ink}\rule{\linewidth}{1.6pt}}\par\vspace{5mm}}
\newcommand{\sectiunemare}[2]{%
  \vspace{4mm}%
  {\noindent\color{ink}\rule{\linewidth}{1.2pt}}\par\vspace{1.5mm}%
  {\noindent\huge\bfseries\color{ink}#1}\par\vspace{1mm}%
  {\noindent\small\itshape\textcolor{grey}{#2}}\par\vspace{1.5mm}%
  {\noindent\color{ink}\rule{\linewidth}{1.2pt}}\par\vspace{3mm}}
\newcommand{\nivel}[3]{%
  \vspace{5mm}%
  \colorlet{lvl}{#1}%
  {\noindent\color{#1}\rule{\linewidth}{2.2pt}}\par\vspace{1.5mm}%
  {\noindent\LARGE\bfseries\color{#1}#2}\par\vspace{1mm}%
  {\noindent\small\itshape\textcolor{grey}{#3}}\par\vspace{3mm}}
\newcommand{\nivelE}[2]{%
  \colorlet{lvl}{#1}%
  \vspace{4mm}{\noindent\color{#1}\rule{\linewidth}{1.5pt}}\par\vspace{1mm}%
  {\noindent\Large\bfseries\color{#1}#2}\par\vspace{2mm}}
\newcommand{\prob}[3]{%
  \vspace{3mm}%
  {\noindent\color{#1}\rule{3pt}{\baselineskip}\hspace{4pt}%
   \textbf{Problema #2}%
   \ifx&#3&\else\ {\small\itshape\textcolor{grey}{--- #3}}\fi}\par\vspace{1mm}}
\newcommand{\sol}{\par\vspace{2mm}{\noindent\small\textcolor{solcol}{\textbf{Soluție.}}\ }}
\newenvironment{parti}{\begin{enumerate}[label=\textbf{(\alph*)},leftmargin=*,itemsep=2pt]}{\end{enumerate}}

%% ---- titluri capitole armonizate ----
\titleformat{\chapter}[block]
  {\normalfont\huge\bfseries\color{ink}}{}
  {0pt}{\vspace*{2mm}\noindent\textcolor{ink}{\rule{\linewidth}{2.4pt}}\vspace{2mm}\newline}
  [\vspace{2mm}\noindent\textcolor{ink}{\rule{\linewidth}{2.4pt}}\vspace{5mm}]
\titleformat{\section}[block]
  {\normalfont\Large\bfseries\color{ink}}{}
  {0pt}{\vspace*{1mm}\noindent\textcolor{ink}{\rule{\linewidth}{1.2pt}}\vspace{1.5mm}\newline}
  [\vspace{1.5mm}\noindent\textcolor{ink}{\rule{\linewidth}{1.2pt}}\vspace{3mm}]
\titleformat{\subsection}[hang]
  {\normalfont\large\bfseries\color{ink}}{}
  {0pt}{}[\vspace{1mm}\noindent\textcolor{grey}{\rule{\linewidth}{0.5pt}}\vspace{2mm}]

%% ---- antet/subsol ----
\setlength{\headheight}{22pt}
\pagestyle{fancy}\fancyhf{}
\lhead{\small\textcolor{grey}{\titluserie}}
\rhead{\small\textcolor{grey}{\capitolcurent}}
\cfoot{\small\textcolor{grey}{\thepage}}
\renewcommand{\headrulewidth}{0.2pt}
\renewcommand{\headrule}{\hbox to\headwidth{\color{grey}\leaders\hrule height \headrulewidth\hfill}}


%% ---- macro-uri suplimentare algebra (Stil B: clasa a XII-a) ----
\colorlet{lvl}{olm}
\newcommand{\Mn}{M_n(\mathbb{C})}
\newcommand{\On}{O_n}
\newcommand{\demo}{\par\vspace{1.6mm}\noindent\textit{\textcolor{solcol}{Demonstrație.}}\ \ignorespaces}
\newcounter{cat}
\newtheoremstyle{probstyle}
  {8pt}{3pt}{\normalfont}{0pt}{\bfseries\large}{}{0.55em}
  {{\color{lvl}\rule[-2.5pt]{3.2pt}{15pt}\hspace{0.55em}}\color{ink}\thmname{#1}\thmnumber{ #2}\thmnote{ \normalfont\small\textcolor{grey}{--- #3}}}
\theoremstyle{probstyle}
\newtheorem{problemaX}{Problema}
\renewcommand{\theproblemaX}{\ifcase\value{cat}\or OLM\or OJM\or ONM\fi.\arabic{problemaX}}
\newenvironment{problema}[1][]
  {\par\vspace{2.5mm}{\color{grey!40}\hrule height 0.4pt}\vspace{0.5mm}\ifx\relax#1\relax\begin{problemaX}\else\begin{problemaX}[#1]\fi}
  {\end{problemaX}}
\newcommand{\theproblema}{\theproblemaX}
\newtheorem*{obsX}{Observație}
\newenvironment{obs}{\begin{obsX}\small\color{ink}}{\end{obsX}}
\newtheorem*{lemat}{Lemă}
\renewcommand{\proofname}{\textcolor{solcol}{Demonstrație}}
\newenvironment{solutie}
  {\par\vspace{1.6mm}\noindent\textit{\textcolor{solcol}{\textbf{Soluție.}}}\enspace\ignorespaces}
  {\hfill\textcolor{solcol}{$\square$}\par\medskip}
\setcounter{tocdepth}{2}
\AtBeginDocument{\renewcommand{\proofname}{Demonstrație}}
\renewcommand{\qedsymbol}{$\blacksquare$}
\setlength{\parskip}{2pt}
\color{ink}

%% ---- parametri de volum (se redefinesc in fisierul fiecarui volum) ----
\newcommand{\titluserie}{Analiză pentru Olimpiadă}
\newcommand{\numarvolum}{1}
\newcommand{\titluvolum}{Analiză, clasa a XI-a}
\newcommand{\subtitluvolum}{}
\newcommand{\isbnvolum}{}
\newcommand{\despreacestvolum}{}
\newcommand{\celelaltevolume}{}

%% ---- compatibilitate tex4ht (build-ul EPUB) ----
%% tex4ht nu poate traduce \halign-ul din mediile *smallmatrix; la conversia in
%% MathML le inlocuim cu variantele normale, care produc <mtable> corect.
\ifdefined\HCode
  \renewenvironment{psmallmatrix}{\begin{pmatrix}}{\end{pmatrix}}
  \renewenvironment{bsmallmatrix}{\begin{bmatrix}}{\end{bmatrix}}
  \renewenvironment{vsmallmatrix}{\begin{vmatrix}}{\end{vmatrix}}
  \renewenvironment{smallmatrix}{\begin{matrix}}{\end{matrix}}
\fi

%% ---- trimiteri catre alt volum (capitolul Constructii e comun volumelor 3 si 4) ----
%% \refv{eticheta}{numar}{volum}: daca eticheta e in volumul curent, \ref obisnuit; altfel numarul ei
%% din volumul in care se afla, urmat de "din volumul N".
\makeatletter
\newcommand{\invol}[1]{\ifnum\numarvolum=#1\relax\else\ din volumul~#1\fi}
\newcommand{\refv}[3]{\@ifundefined{r@#1}{#2\invol{#3}}{\ref{#1}}}
\makeatother
